%%%%%%%%%%%%%%%%%%%%%%%%%%%%%%%%%%%%%%%%%%%%%%%%%%%%%%%%%%%%%%%%%%%%%%%%
\section{GDL-Bench Challenge Problem (20 points)}
%%%%%%%%%%%%%%%%%%%%%%%%%%%%%%%%%%%%%%%%%%%%%%%%%%%%%%%%%%%%%%%%%%%%%%%%

\begin{question}
As part of each homework, you will contribute one candidate problem to GDL-Bench, a benchmark for evaluating the reasoning capabilities of current AI systems in generative and geometric deep learning. For this homework, design one original challenge problem related to one or more of the topics covered in this assignment. The problem should require genuine mathematical, conceptual, algorithmic, or implementation-level reasoning, and should not be answerable by simply recalling a definition or looking up a fact.
\begin{enumerate}
  \item[(a)] \textbf{Challenge problem.} Write the exact problem as you would like it to appear in GDL-Bench. It must be original, self-contained, technically well-defined, and unambiguous.
  \item[(b)] \textbf{Gold answer and solution.} Provide the correct final answer together with a concise but complete solution, enough for another student or instructor to independently verify the answer. Whenever possible, design the problem so that the final answer can be objectively checked.
  \item[(c)] \textbf{What does the problem test?} Topic; problem type (derivation, conceptual reasoning, counterexample, numerical, coding/debugging, algorithm design, or other); and the core insight in 1--3 sentences, including why a solver with only superficial knowledge might struggle.
  \item[(d)] \textbf{Test AI systems.} Only after the problem and solution are finalized, give the exact same problem to at least two current AI models, with no extra hints and without the solution. For each model report the name/version, whether the answer was Correct, Partially Correct, or Incorrect, and a 1--3 sentence explanation of any important error. A problem does not need to defeat the models to receive full credit.
\end{enumerate}
Quality bar: original, relevant, self-contained, unambiguous, verifiable, and reasoning-intensive. Grading: relevance and originality (4); technical correctness (6); clarity (4); depth of reasoning (4); complete metadata and AI evaluation (2).
\end{question}

The full, agent-runnable version of this problem (task statement, reference
solution, and automated tests) is in \texttt{gdlbench\_submission/}. The
statement below is the text of \texttt{task/TASK.md}.

\subsection{(a) Challenge problem}

\textbf{Which generator does each GAN objective prefer?}

All logarithms are natural, and all distributions live on the nonnegative
integers. The data distribution \(p\) is uniform on \(\{0,1,2,3\}\). For each
positive integer \(k\), the generator distribution \(q_k\) is uniform on
\(\{0,1,\dots,k-1\}\). Conventions: \(\log 0=-\infty\),
\(0\cdot(\pm\infty)=0\), and an expectation that puts positive probability on
an integrand equal to \(+\infty\) is \(+\infty\).

\medskip\noindent\textbf{Definitions.}
\begin{itemize}
  \item \(\mathrm{JSD}(p,q)=\tfrac12\KL(p\|m)+\tfrac12\KL(q\|m)\), with \(m=(p+q)/2\).
  \item \emph{Minimax GAN.} For a discriminator \(D\) mapping integers to \([0,1]\),
  \[
    V(D,q)=\E_{x\sim p}[\log D(x)]+\E_{x\sim q}[\log(1-D(x))].
  \]
  Let \(D^*_q\) maximize \(V(D,q)\) (its values off the support of \(p+q\) do
  not matter), and define
  \[
    V_\star(q)=V(D^*_q,q),\qquad
    L_{\mathrm{MM}}(q)=\E_{x\sim q}[\log(1-D^*_q(x))],\qquad
    L_{\mathrm{NS}}(q)=\E_{x\sim q}[-\log D^*_q(x)].
  \]
  \item \emph{Least-squares GAN} with labels \(a=-1\) (generated), \(b=1\)
  (real), and generator target \(c=0\). The discriminator \(D^{\mathrm{LS}}_q\)
  minimizes \(\tfrac12\E_p[(D-1)^2]+\tfrac12\E_q[(D+1)^2]\) over real-valued
  \(D\), and the generator criterion is
  \[
    V_{\mathrm{LS}}(q)=\tfrac12\E_{x\sim p}\bigl[(D^{\mathrm{LS}}_q(x)-0)^2\bigr]
    +\tfrac12\E_{x\sim q}\bigl[(D^{\mathrm{LS}}_q(x)-0)^2\bigr].
  \]
\end{itemize}
The generator wants each of \(V_\star\), \(L_{\mathrm{MM}}\),
\(L_{\mathrm{NS}}\), and \(V_{\mathrm{LS}}\) to be as small as possible.

\medskip\noindent\textbf{Report:}
\begin{enumerate}
  \item \texttt{jsd\_tie\_pairs}: every unordered pair \(\{j,k\}\) of distinct
        positive integers with \(\mathrm{JSD}(p,q_j)=\mathrm{JSD}(p,q_k)\) exactly.
  \item \texttt{ls\_tie\_pairs}: every unordered pair \(\{j,k\}\) of distinct
        positive integers with \(V_{\mathrm{LS}}(q_j)=V_{\mathrm{LS}}(q_k)\) exactly.
  \item \texttt{ls\_value\_q2}: the value of \(V_{\mathrm{LS}}(q_2)\).
  \item \texttt{k\_hat}: the integer \(\hat k\neq2\) such that \(\{2,\hat k\}\)
        is in \texttt{jsd\_tie\_pairs}, or \texttt{null} if there is none.
  \item If \(\hat k\) exists: the values of \(V_\star\), \(L_{\mathrm{MM}}\),
        and \(L_{\mathrm{NS}}\) at \(q_2\) and at \(q_{\hat k}\).
  \item If \(\hat k\) exists: for each of \(V_\star\), \(L_{\mathrm{MM}}\),
        \(L_{\mathrm{NS}}\), which of \(q_2\) and \(q_{\hat k}\) the generator
        prefers (lower value), or \texttt{"tie"}.
\end{enumerate}
The solver also writes a short derivation that justifies why the lists in
items 1 and 2 are \emph{complete}. Answers go in
\texttt{task/results/answer.json}, with \(+\infty\) written as the string
\texttt{"inf"} and floats to at least 10 significant figures. Pair lists,
\(\hat k\), \texttt{"inf"}, and preferences must match exactly; floats must be
within \(10^{-8}\).

\subsection{(b) Gold answer and solution}

\begin{center}
\begin{tabular}{@{}lll@{}}
\toprule
Key & Gold answer & Value\\
\midrule
\texttt{jsd\_tie\_pairs} & \(\{1,16\},\ \{2,8\}\) & \\
\texttt{ls\_tie\_pairs} & \(\{1,16\},\ \{2,8\}\) & \\
\texttt{ls\_value\_q2} & \(1/3\) & \(0.333333\)\\
\texttt{k\_hat} & \(8\) & \\
\(V_\star(q_2)=V_\star(q_8)\) & \(\log2-\tfrac32\log3\) & \(-0.954771\)\\
\(L_{\mathrm{MM}}(q_2)\) & \(\log(2/3)\) & \(-0.405465\)\\
\(L_{\mathrm{MM}}(q_8)\) & \(-\tfrac12\log3\) & \(-0.549306\)\\
\(L_{\mathrm{NS}}(q_2)\) & \(\log3\) & \(1.098612\)\\
\(L_{\mathrm{NS}}(q_8)\) & \(+\infty\) & \\
\texttt{prefers} & \multicolumn{2}{l}{\(V_\star\): tie;\quad \(L_{\mathrm{MM}}\): \(q_8\);\quad \(L_{\mathrm{NS}}\): \(q_2\)}\\
\bottomrule
\end{tabular}
\end{center}

\paragraph{Both criteria depend only on a support ratio.} The distributions
\(p\) and \(q_k\) are nested uniforms. Let the smaller one have support size
\(n\) and the larger one \(N\), and let \(r=N/n=\max(k,4)/\min(k,4)\ge1\). On
the shared support \(m=(r+1)/(2N)\); off it, \(m=1/(2N)\). Then
\[
  \KL(\text{small}\|m)=\log\frac{2r}{r+1},\qquad
  \KL(\text{large}\|m)=\frac1r\log\frac{2}{r+1}+\Bigl(1-\frac1r\Bigr)\log2,
\]
so
\[
  \mathrm{JSD}(r)=\log2+\tfrac12\log r-\tfrac12\Bigl(1+\tfrac1r\Bigr)\log(1+r),
  \qquad
  \frac{d\,\mathrm{JSD}}{dr}=\frac{\log(1+r)}{2r^2}>0 .
\]
JSD is symmetric, so it does not matter which of \(p\) and \(q_k\) is the
larger one. JSD is therefore a strictly increasing function of \(r\) alone,
and \(\mathrm{JSD}(p,q_j)=\mathrm{JSD}(p,q_k)\) holds iff \(r_j=r_k\). For
\(j\neq k\) this needs one index below 4 and one above it, with \(4/j=k/4\),
i.e.\ \(jk=16\). The only integer solutions are \(\{1,16\}\) and \(\{2,8\}\)
(\(j=3\) gives \(16/3\)). The argument covers every positive integer, so the
list is complete, and \(\hat k=8\).

\paragraph{LSGAN.} Minimizing \(\tfrac12p(D-1)^2+\tfrac12q(D+1)^2\) pointwise
gives \(D^{\mathrm{LS}}=(p-q)/(p+q)\). Substituting,
\[
  V_{\mathrm{LS}}=\frac12\sum_x\frac{(p-q)^2}{p+q}
  =\frac12\Bigl[\frac{(r-1)^2}{r(r+1)}+\frac{r-1}{r}\Bigr]
  =\frac{r-1}{r+1},
\]
which is also strictly increasing in \(r\) alone. The tie pairs are the same,
and \(V_{\mathrm{LS}}(q_2)=1/3\).

\paragraph{Why the ties happen.} An \(f\)-divergence depends only on the
distribution of the likelihood ratio \(p/q\). The pair \((p,q_2)\) has the
same ratio structure as \((q_8,p)\): a uniform against a uniform on half its
support. Hence \(D_f(p\|q_2)=D_g(p\|q_8)\) with the reversed generator
\(g(u)=u\,f(1/u)\). Every symmetric \(f\)-divergence (JSD, the LSGAN
\(\chi^2\)-type term) must tie on this pair. Only asymmetric ones can tell
mode dropping (\(q_2\)) apart from over-spreading (\(q_8\)).

\paragraph{GAN terms.} \(D^*=p/(p+q)\).
\begin{itemize}
  \item \(q_2\): \(D^*=1/3\) on \(\{0,1\}\), \(D^*=1\) on \(\{2,3\}\).
  Then \(V_\star=\tfrac12\log\tfrac13+\log\tfrac23=\log2-\tfrac32\log3\),
  \(L_{\mathrm{MM}}=\log\tfrac23\), and \(L_{\mathrm{NS}}=\log3\).
  \item \(q_8\): \(D^*=2/3\) on \(\{0,\dots,3\}\), \(D^*=0\) on \(\{4,\dots,7\}\).
  Then \(V_\star=\log\tfrac23+\tfrac12\log\tfrac13+\tfrac12\log1=\log2-\tfrac32\log3\),
  \(L_{\mathrm{MM}}=\tfrac12\log\tfrac13=-\tfrac12\log3\), and
  \(L_{\mathrm{NS}}=\tfrac12\log\tfrac32+\tfrac12(+\infty)=+\infty\).
\end{itemize}
\(V_\star=2\,\mathrm{JSD}-\log4\) ties, as it must. The generator's own minimax
term does \emph{not} tie: the \(\E_p[\log D^*]\) term also depends on \(q\)
through \(D^*\), and dropping it leaves a quantity that favors \(q_8\). The
non-saturating loss is infinite for any generator that puts mass where the
data has none, so it favors \(q_2\), the generator that drops modes. The same
game value gives three different verdicts.

\paragraph{Verification.} \texttt{gt\_code/reference\_solution.py} evaluates
every quantity directly from the distributions (exact fractions wherever no
logarithm is involved), scans \(k=1,\dots,512\) for ties, and asserts
agreement with the closed forms and with \(V_\star=2\,\mathrm{JSD}-\log4\).
\texttt{tests/test\_solution.py} passes 8/8 on the gold answer, and each of
the following mistakes fails at least one test:
\begin{center}
\begin{tabular}{@{}ll@{}}
\toprule
Mistake & Fails\\
\midrule
Report only \(\{2,8\}\) (miss the pair below 4) & \texttt{test\_jsd\_tie\_pairs}\\
Treat \(L_{\mathrm{MM}}\) like \(V_\star\) (claim it ties) & \texttt{test\_q\_k\_hat\_terms}, \texttt{test\_prefers}\\
Assume non-saturating has the same optimum (``tie'') & \texttt{test\_prefers}\\
Drop the \(\tfrac12\) in \(V_{\mathrm{LS}}\) (report \(2/3\)) & \texttt{test\_ls\_value\_q2}\\
Finite \(L_{\mathrm{NS}}(q_8)\) (ignore \(D^*=0\) off the data support) & \texttt{test\_q\_k\_hat\_terms}\\
\bottomrule
\end{tabular}
\end{center}

\subsection{(c) What does the problem test?}

\begin{itemize}
  \item \textbf{Topic:} GAN objectives and the divergences they minimize
        (minimax/JSD, non-saturating loss, LSGAN/Pearson-type \(\chi^2\)),
        mode dropping vs.\ mode covering.
  \item \textbf{Problem type:} derivation, conceptual reasoning, and
        counterexample.
  \item \textbf{Core insight:} symmetric \(f\)-divergences such as JSD and
        the LSGAN criterion depend only on the support ratio, so they cannot
        distinguish dropping half the modes from spreading over twice the
        support; the map \(k\mapsto16/k\) produces exactly the ties. A solver
        with only surface knowledge repeats that ``GANs penalize mode
        collapse'' and treats the game value, the minimax generator term, and
        the non-saturating loss as interchangeable, when in fact they give
        three different preferences on the same pair.
\end{itemize}

\subsection{(d) Testing AI systems}

Each model ran as a coding agent in Cursor, given only \texttt{task/} in a
fresh workspace and the same fixed prompt (no hints, no internet lookup of
the problem, no access to tests or solutions). Each \texttt{answer.json} was
graded with \texttt{tests/test\_solution.py}.

\begin{center}
\begin{tabular}{@{}lll@{}}
\toprule
Model (version) & Result & Tests passed\\
\midrule
Claude Opus 5.5 & Correct & 8/8\\
GPT-5.6-sol & Correct & 8/8\\
Grok 4.7 & Correct & 8/8\\
\bottomrule
\end{tabular}
\end{center}

\begin{itemize}
  \item \textbf{Claude Opus 5.5:} correct. It showed both criteria depend only
        on the support ratio and are monotone in it, which reduces ties to
        \(jk=16\). It also noted \(V_\star=2\,\mathrm{JSD}-\log4\) and
        cross-checked by scanning \(k\le400\) with exact fractions.
  \item \textbf{GPT-5.6-sol:} correct. It gave the same reduction to
        \(u=\min(k/4,4/k)\) with an explicit derivative, and the cleanest
        written derivation of the three.
  \item \textbf{Grok 4.7:} correct, and the most thorough. It derived
        separate closed forms on each side of \(k=4\), proved the identity
        \(\varphi(16/j)=\mathrm{JSD}(p,q_j)\) explicitly, and compared
        values exactly over \(k\le16\).
\end{itemize}
All three found the \(\{1,16\}\) pair, proved the tie lists complete, and
separated \(V_\star\), \(L_{\mathrm{MM}}\), and \(L_{\mathrm{NS}}\) correctly.
None fell into the anticipated traps. Current agentic models with code
execution handle closed-form divergence calculations reliably, so this
problem works best as a clean, verifiable check of the symmetric-divergence
insight and of the difference between the three generator objectives, rather
than as an item that breaks models.
